\documentclass[a6paper,fontsize=10pt,twoside,openright]{scrbook}
\usepackage[a6paper,top=1cm,bottom=1.9cm,left=1cm,right=1cm,footskip=1cm,headsep=0.2cm]{geometry}% add showframe
\usepackage[swedish]{babel}
\usepackage{graphicx,tocloft,metalogo,titletoc,fontspec,titlesec,eso-pic,scrpage2,anyfontsize,microtype,array,marvosym,mdframed,textcomp,amssymb,enumitem}
\usepackage[columns=1]{idxlayout}

\newgeometry{top=1cm,bottom=1cm,left=1cm,right=1cm,footskip=0pt,headsep=0.2cm}


%\setlength{\emergencystretch}{2em}

\usepackage{letltxmacro}

\LetLtxMacro\origttfamily\ttfamily
\DeclareRobustCommand*{\ttfamily}{%
  \origttfamily
  \hyphenchar\font=`\-\relax
  \fontdimen3\font=.25em\relax
  \fontdimen4\font=.167em\relax
  \fontdimen7\font=.167em\relax
}

%% Font: Baskervald X
%% \usepackage[lf]{Baskervaldx} % lining figures
%% \usepackage[bigdelims,vvarbb]{newtxmath} % math italic letters from Nimbus Roman
%% \usepackage[cal=boondoxo]{mathalfa} % mathcal from STIX, unslanted a bit
%% \renewcommand*\oldstylenums[1]{\textosf{#1}}

%\usepackage[ebgaramond-maths]{mathdesign}
\usepackage[bitstream-charter]{mathdesign}
\usepackage[htt]{hyphenat}

%% Gömma varningar om underfull \hbox
%% \hbadness=10001

%% \tolerance=1
%% \emergencystretch=\maxdimen
%% \hyphenpenalty
%% \exhyphenpenalty


\usepackage{setspace}
\renewcommand{\baselinestretch}{0.98} 

%\usepackage{showframe}
\usepackage{makeidx}
\makeindex

\let\oldemptyset\emptyset
\let\emptyset\varnothing

%% eso-pic for background on chapter pages %%
\newcommand\BackgroundPic{\includegraphics[width=\paperwidth,height=\paperheight,keepaspectratio]{elements/chapter2.pdf}}

%% scrpage2 pagestyling %%
\pagestyle{scrheadings}
\clearscrheadfoot

%% change empty page produced by \cleardoublepage into scrheading style %%
\makeatletter
\let\ps@empty\ps@scrheadings
\makeatother

\newcommand{\rpt}{
        \raisebox{.2ex}{:}
        \hspace*{-0.5ex}\raisebox{-.4ex}{\rule{.1ex}{2.3ex}\,\rule{.2ex}{2.3ex}}}

\newcommand{\revrpt}{
        \raisebox{-.4ex}{\rule{.2ex}{2.3ex}\,\rule{.1ex}{2.3ex}}
        \hspace*{-0.5ex}\raisebox{.2ex}{: }}

%% chapter and section styling %%
\titleformat{\chapter}[hang]{}{}{0pt}{\centering\fontsize{30}{70}\rmfamily\textbf}
\titleformat{\section}[hang]{}{}{0pt}{\fontsize{13}{15}\rmfamily\textbf}
\titlespacing*{\section}{0pt}{0pt}{0pt}
\titlecontents{chapter}[0em]{\addvspace{3pt}}{\contentslabel{}}{}
              {\titlerule*[0.3pc]{.}\hspace*{-4pt}\contentspage}

%\newfontfamily\chapterfont[]{ChopinScript}

%\addtokomafont{pagenumber}{\oldstylenums}
              
%% Command for chapter page, will remove foot and head %%
\newcommand{\chapterpage}[2]{
  \newpage
  \null
  \cleardoublepage
  %\AddToShipoutPicture*{\BackgroundPic}
  \chapter[#1]{#2}
  \renewcommand{\leftmark}{#1}
  \newpage
  \ohead{\textnormal{\textsc{\scriptsize\leftmark}}}
  \ofoot[\pagemark]{\textsc{\scriptsize\pagemark}}
  %\cfoot{\centering\includegraphics[keepaspectratio]{elements/footelement2.pdf}}
}
\newcommand{\hiddenchapterpage}[1]{
  \newpage
  \null
  \cleardoublepage
  %\AddToShipoutPicture*{\BackgroundPic}
  \chapter*{#1}
  \renewcommand{\leftmark}{#1}
  \newpage
  \ohead{\textnormal{\textsc{\scriptsize\leftmark}}}
  \ofoot[\pagemark]{\textsc{\scriptsize\pagemark}}
  %\cfoot{\centering\includegraphics[keepaspectratio]{elements/footelement2.pdf}}
}
\makeatletter
\renewcommand\chapter{\if@openright\cleardoublepage\else\clearpage\fi
  \thispagestyle{empty}
  \global\@topnum\z@
  \@afterindentfalse
  \ofoot[\pagemark]{\textsc{\scriptsize\pagemark}}
  \cfoot{}
  \ohead{}
  \secdef\@chapter\@schapter}
\makeatother

\newcommand{\specialcell}[2][c]{
  \begin{tabular}[#1]{@{}l@{}}#2\end{tabular}} 

%% Command for songs %%
\newcommand{\visa}[1]{
  \index{#1}
  \section{#1}
}
\newcommand{\fvisa}[2]{
  \index{#1}\index{#2}
  \section{#1}
}
\newcommand{\svisa}[3]{
  \index{#1}\index{#2}
  \section{#3}
}

%% TOC %%
\setcounter{tocdepth}{0}
\tocloftpagestyle{scrheadings}

%% MISC %%
\pagenumbering{arabic}
\renewcommand{\chaptername}{}
\renewcommand{\thechapter}{}
\renewcommand{\thesection}{}

\makeatletter
\renewenvironment{theindex}
                 {{\noindent\rmfamily{\fontsize{13}{15}\textbf{TITEL och FÖRSTARADEN REGISTER}}\vspace*{10pt}\par\textbf{\normalsize \#}\par}
                   \@mkboth{\MakeUppercase\indexname}
                           {\MakeUppercase\indexname}
                           \let\item\@idxitem}
                 {}
\makeatother
\newmdenv[
  topline=false,
  bottomline=false,
  rightline=false,
  linewidth=1pt,
  skipabove=5pt,
  skipbelow=3pt,
  innertopmargin=0pt,
  innerbottommargin=0pt
]{leftborder}
\makeatletter

\newcommand*\cleartoleftpage{%
  \clearpage
  \ifodd\value{page}\hbox{}\newpage\fi
}

\begin{document}
\setlength{\parindent}{15pt}
\vspace*{6.5cm}
\hspace*{0.9cm}
\includegraphics[keepaspectratio,width=7cm]{elements/logo.pdf}
\clearpage
\noindent Denna förträffliga oersättliga Manual tillhör:
\ohead{\textnormal{\textsc{\scriptsize\leftmark}}}
\ofoot[\pagemark]{\textsc{\scriptsize\pagemark}}
\clearpage
%\section{Datavetare är skapande människor}\vspace{10pt}
\noindent{\ttfamily Datavetare är skapande människor\\ \\ Datavetare
  skapar. De skapar nya program som världen aldrig sett tidigare, för
  att låta datorer göra saker som de aldrig gjort tidigare. De forskar
  fram nya kunskaper och öppnar nya fält för framtidens
  forskning. Datavetare skapar de verktyg som gör det möjligt att lösa
  nya problem. Nya programmeringsspråk, nya beräkningsverktyg, bättre
  kompilatorer, bättre databaser, nya operativsystem, bättre
  utvecklingsverktyg, distribuerade system... Listan kan göras hur
  lång som helst.\\ \\ Datavetenskapens handfasta mål är att förbättra
  den programvara som vi kör på våra datorer. Vad som krävs är att
  programmen gör vad de skall (korrekthet), att de inte gör fel om
  något oförutsett inträffar (pålitlighet), och att de arbetar
  effektivt (effektivitet). Dessutom skall programmen utvecklas på ett
  effektivt sett.}\\
\\
{\footnotesize\textit{Ur ``Vad är en datavetare egentligen''\\
    BlurGel nr 0, 1995}}
\clearpage
%\newgeometry{top=1cm,bottom=1.2cm,left=1cm,right=1cm,footskip=0pt,headsep=0.2cm}
\null
\vfill
    {\noindent\small\centering
      Redaktör: ?\\
      Grafisk formgivare: Per Bergqwist\\
      Logotyp av Jens Bohlin\\
      Producerad med \XeTeX\\
      v1.0: 1994-11-21 (1000 exemplar)\\
      v1.1: 1996-01-02 (1000 exemplar)\\
      v2.001: 2001-11-15 (1000 exemplar)\\
      v11111011111: 2015-??-?? (1200 exemplar)\\
      Tryckeri: Boktryckarna\par
    }
%\restoregeometry
\cleardoublepage
%\cfoot{\centering\includegraphics[keepaspectratio]{elements/footelement2.pdf}}
\section{FÖRORD}\vspace{10pt}
%\noindent{\large\textrm{\textbf{FÖRORD}}\vspace*{10pt}}\\
\hspace{10pt} ost
\cleardoublepage
\section{FÖRORD från v2.001}\vspace{10pt}
%\noindent{\large\textrm{\textbf{FÖRORD från v2.001}}\vspace*{10pt}}\\
\hspace{10pt}Allt sedan Manualen v1.1 kom till världen har det
funnits en skara människor som ägnat sina liv åt att vänta på dess
efterföljare. Denna skaras oförtrutliga tillväxt utmynnande i att ett
formaliserat uppdateringsarbete satte igång våren 1999. Den 4 december
-99 samlades en liten men representativ skara från C i Linköping, DVP
i Uppsala och C i Umeå för att samla alla tankar och idéer som
framförts på den då flitigt utnyttjade mailinglistan
\texttt{manual@acc.umu.se}. Efter detta möte trodde vi att arbetet
skulle leda till en färdig Manual våren 2000. Så blev det dock inte på
grund av en utdragen och högst osaklig debatt om en del låtar som
redan var bortplockade. Nu, över ett år senare, är vi som har
funderat, fipplat och till och med slitit(?) med denna nya bok
lyckliga över att få presentera Manualen version 2.001\\ \indent En
betydande samling sånger som Manualen kräver naturligtvis en avsevärd
mängd arbete att sammanställa, faktum är att antalet välvilliga
medhjälpare från alla berörda order förmodligen inte går att passa in
i en nibble. Det stora antalet inblandade parter gör att en uppräkning
är dömd att innehålla felaktigheter, dock gör vi nedan ett dj-ärvt
försök samtidigt som vi hoppas att alla vi glömt kan förlåta oss i
någon framtid.
\newpage
\indent Vilka var de tappra själarna som var med på mötet i Uppsala
då?\\
\\
\noindent
\begin{tabular}{@{}p{0.5\textwidth}p{0.5\textwidth}@{}}
  Lars-Arne Mattson & Anders Ramsell\\
  Henrik Eriksson & Olof Dahlberg\\
  Hans Molin & Björn Bergman\\
  Jesper Wilhelmsson & Ingela Anderton\\
  Daniel Gref & Mattias Hallberg\\
  Daniel Holmgren & Andreas Hed
\end{tabular}\\
\\
%% \\
%% Manualen finns också på nätet,\\ URL:
%% \texttt{http://www.acc.umu.se/\raisebox{.4ex}\texttildelow
%%   manualen/}\\
\\ Hoppas nu att ni är nöjda med
resultatet.\\ \\ Uppsala den 24 december 2000 (God Jul!)\\ \\ Björn
Bergman / Jesper Wilhelmsson / Andreas Hed\\ \\ De sista justeringarna
skedde under den datavetenskapliga utbildningskonferensen i Umeå 2001.
\cleardoublepage
\section{FÖRORD från v1.0}\vspace{10pt}
%\noindent{\large\textrm{\textbf{FÖRORD från v1.0}}\vspace*{10pt}}\\
\hspace{10pt}Detta är en gemensam sångbok för samtliga datavetare
i Sverige; C i Lindköping, DVP i Uppsala och C i Umeå. Vi har genom
denna bok fått ytterligare en sak gemensamt.\\ \indent Att skapa en
gemensam sångbok har diskuterats flerga gånger, jag minns att
kontakter tog under 100-köret 1990 eller var det 1991?  Saker som
händer under 100-kör är inget man minns, utan saker som man senare
tror sig har hänt?! Den gången brakade arbetet samman efter några
veckor. Men idén levde kvar och gav 3 (eller 4) år senare detta som
resultat.\\ \indent Sångboken bygger på ``DVL:s Lilla blå v3.0\ss''
som aldrig han bli tryckt, innan arbetet kom igång för att skapa denna
bok. Dock skall sägas att denna bok är betydligt större än ``Lilla blå
v3.0\ss'', omkring 200 sidor för att inte vara exakt.\\ \indent När
\textit{manualen} trycks i december 1994 är boken den största sångbok
som finns runt om i vårt rike! Ok Psalmboken är större, men inte lika
roligt att ha på gasquer!\\ \indent Vi som varit ansvariga för denna
bok, har försökt att i så stor utsträckning som möjligt att få med
melodier, låtarnas ursprung, och hur man skall sjunga. Allt för att
underlätta för folk som inte är i kontakt med lika mycket sånger som
vissa av oss andra.\\ \indent Det är många som har hjälpt till med
denna bok, inte enbart datavetare utan även personer från helt andra
linjer och orter, jag vill rikta ett speciellt tack till Er.
\newpage
\indent Manualens generalgenomsjungning ägde rum på\\ Strålsnäs
herrgård den 18 November 1994, då vi avnjöt en underbar
rådjursmiddag. På denna sittning sjöng vi så mycket vi orkade och hade
efter ca 4 timmar traverserat ca 40 sånger.\\ \indent Med på denna
sångboksgenomsjungning och därav lika (o)ansvariga som jag till de
eventuella felen var,\\ \textit{Micke Liljeblad, Mats Hägglund, Jakob
  Engblom, Pär\\ Mattsson, Martin Davidson, Anders Jackson,
  Björn\\ Knutsson, Måns Engmem, Måns af Klercker, Jonas\\ Högström,
  Calle Englund, Andreas Gref} och sist jag själv.\\ \indent Innan jag
kommer till alla som skall tackas vill jag bara nämna att en sångbok
är till för att sjungas ur, detta betyder inte att man behöver
sjunga. Följande personer är jag evigt(?)  tacksam till;\\

\noindent
\begin{tabular}{@{}p{0.5\textwidth}p{0.5\textwidth}@{}}
  Anders Jansson & Anders Wilhelm\\
  Bertil Spolander & Björn Knutsson\\
  Björn Torkelsson & Erik Åhlin\\
  Göran Svensson & Håkan Lindholm\\
  Ingela Ljungkvist & Jan-Olof Flink\\
  Jens Bohlin & Joakim Thörnkvist\\
  Jonas Högström & Lars Hägglund\\
  Måns af Klercker & Olof Appelblad\\
  Pär Mattson
\end{tabular}
\newpage
\indent Dessutom vill jag tacka alla som lusläst boken, sammt alla Ni
som gjorde förhandsbeställningar. Det finns säkert många andra som jag
missat och jag är precis lika tacksam till er som alla
andra!\\
%% \indent Om ni har några bidrag eller saker som kan vara värda
%% att rätta till så skicka dessa till någon av festerierna nedan;\\

%% \noindent
%% \begin{tabular}{@{}p{0.3\textwidth}p{0.7\textwidth}@{}}
%%   Piraya klubben & piraya@ts.umu.se\\
%%   Klubbverket & dvl-klubbverket@student.docs.uu.se\\
%%   CC & cc@cyd.liu.se
%% \end{tabular}\\

%% \indent Till sist, kom ihåg att hela sångboken och fellista finns på:
%% \texttt{http://www.ts.umu.se/\raisebox{.4ex}\texttildelow manualen/}\\
%% \\
\noindent Umeå den 21 November
1994\\ \\ NeNNe
\cleartoleftpage
\section{TRAVERSERING}\vspace{10pt}
\hspace{10pt}Genom hela denna bok gäller följande traverseringsregler. I sånger där
en försångare eller dylikt skall sjunga ensam används fet text. Ifall
något speciellt såsom en handling inträffar vid ett ord, så används
(kursiv), där ordet för handling eller beteende står inom ( ). En
sådan handling gäller oftast maximalt raden ut. Ifall en längre
förklaring behövs kursiveras texten. Beskrivningen finns sedan efter
sångens slut. I denna beskrivning kan mycket annat matnyttigt finnas,
läs den gärna.\\ \indent I vissa sånger skall könen sjunga olika
rader, i dessa sånger används {\Large\Male} för manligt kön och där
kvinnligt kön skall sjunga används givetvis då {\Large\Female}.
\begin{leftborder}
  \hspace{10pt}Då det ibland kan förekomma lokala variantioner på vissa sånger, så
  används en  vertikal linje i högra kanten för att markera att detta
  är olika alternativ. Ibland betyder linjen att det är raderna som är
  alternativen, och ibland betyder linjen att det är olika
  versalternativ. Men den kloke läsaren förstår säkert när vad är vad.
\end{leftborder}
\hspace{15pt} För att göra allt ännu svårare så skall man oftast sjunga med stark
stämma där VERSALER används, men observera att man inte alltid skall
göra detta. Vem har sagt att det skall vara lätt?
\newpage
%% begin toc on odd page and style the title %%
\cleardoublepage
\renewcommand{\contentsname}{\vspace{-2.17cm}\rmfamily{\fontsize{13}{15}\textbf{INNEHÅLL}}\vspace{-1.2cm}}
\tableofcontents

%% Remove the toc title from heading %%
\renewcommand{\leftmark}{}

\setlength{\parindent}{0pt}
\raggedbottom
\chapterpage{Av tradition}{Av tradition}
\null
\vfill
\hfill\begin{minipage}[c]{6.7cm}
\textit{Vad bjuder oss uppriktigt Afrika?\\
Vad visa kan Amerika?\\
Vad Asien? Vad allt Europa?\\
Jag trotsar öppet alltihopa.\\
Men Skandinavien - det är alladar!\\
Blott Sverige svenska krusbär har.\\
\\
Carl Love Almquist}
\end{minipage}\hfill
\vfill
\null
\newpage
\visa{DU GAMLA, DU FRIA}
Du gamla, Du fria, Du fjällhöga nord\\
Du tysta, Du glädjerika sköna!\\
Jag hälsar Dig, vänaste land uppå jord,\\
||: Din sol, Din himmel, Dina ängder gröna. :||\\
\\
Du tronar på minnen från fornstora dar,\\
då ärat Ditt namn flög över jorden.\\
Jag vet att Du är och Du blir vad du var.\\
||: Ja, jag vill leva jag vill dö i Norden. :||\\
\\
Jag städs vill dig tjäna mitt älskade land,\\
din trohet till döden vill jag svära.\\
Din rätt, skall jag värna, med håg och med hand,\\
||: din fana, högt den bragderika bära. :||\\
\\
Med Gud skall jag kämpa, för hem och för härd,\\
för Sverige, den kära fosterjorden.\\
Jag byter Dig ej, mot allt i en värld\\
||: Nej, jag vill leva jag vill dö i Norden. :||\\
\\
{\footnotesize\textit{Text: Richard Dybeck, 1844\\
\\
Sveriges nationalsång av tradition sedan 1866.\\
De två sista verserna sjungs sällan.}}
\clearpage

\newpage
%\fvisa{SVERIGES FLAGGA}{Flamma stolt mot dunkla skyar}
\vspace{10pt}
Flamma stolt mot dunkla skyar\\
lik en glimt av sommarens sol!\\
Över Sveriges skogar, berg och byar,\\
över vatten och viol!\\
Du som sjunger, när Du bredes\\
som vår gamla lyckas tolk.\\
Solen lyser! Solen lyser!\\
Ingen vredes åska slog vårt tappra folk!\par
\vspace{10pt}
Flamma högt vårt kärlekstecken!\\
Värm oss, när det blåser kallt!\\
Stråla ut de blåa vecken\\
kärlek mera stark än allt!\\
Sveriges flagga! Sveriges ära!\\
Fornklenod och framtidstolk!\\
Gud är med oss! Gud är med oss!\\
Han skall bära stark vårt fria svenska folk\par
\vspace{10pt}
{\footnotesize\textit{Text: K.G. Ossiannilsson\\
Musik: Hugo Alfvén}}

%\fvisa{JÄMTLANDSSÅNGEN}{Så tåga vi tillsammans bort}
\vspace{10pt}
Så tåga vi tillsammans bort\\
mellan Jämtlands gröna ängar\\
bort mellan nyland som prunka\\
fulla av bröllopsblomsters prakt.\\
Så skådom nu med gamman hän\\
emot berg i blåa fjärran hän,\\
över sjöar, strömmar, skogar\\
jämt kring bygder på vakt.\\
\\
\revrpt Fagert är landet\\
som blev vår lott och arvedel.\\
Så firom dess fägring nu\\
med sång och stråkars spel.\\
Så tändom ånyo det hopp\\
som våra fäder närt.\\
För slit och mödor\\
av fröjd och sol\\
ett mått oss beskärt\rpt\\
\\
{\footnotesize\textit{Text \& Musik: Wilhelm Peterson-Berger}}
\clearpage

%\input{songs/jamtlandssången.tex}
\input{songs/studentsången.tex}
\newpage
\visa{UNDER SVEA BANÈR}
\vspace{10pt}
Under Svea Banér Himlen seger oss ger;\\
Då för Konung och Land Äran lyfter sin hand.\\
Ännu Svearnes mod sif bereder en Stod\\
Utaf Lagrarne höjd, fast besglad mes blod.\\
Himlen gifve oss frid!\\
Men om den icke vinne utan vapen och strid,\\
Blifve Segern då vår, ögat offre sin tår\\
Den som faller, ock rätta till vår tacksamhet får.\\
\\
{\footnotesize\textit{Text: Samuel Ödmann\\
Musik: Johann Christian Friedrich Haeffner\\
\\
Anses vara den första upsaloensiska studentsången. Musiken hade
skrivits tidigare som soldatkör till Hæffners opera “Renaud” (\oldstylenums{1801}).\\
\\
Framfördes första gången, mitt under kriget mot Ryssland,
den \oldstylenums24} oktober \oldstylenums{1808}. Detta skedde vid en
uppvakting för fältmarskalkten Klingspor på genomresa från Finland
till Stockholm som gästade landshövding Uppsala slott.  Uppsalas
studenter marscherade då upp till slottet. Sångerna gick i täten och
sjöng denna sång.}}

\newpage
\visa{O, GAMLA KLANG- OCH JUBELTID}
{\footnotesize\textit{Melodi: Oh alte Burschenherrlichkeit}}\par
\vspace{10pt}
O gamla klang och jubeltid\\
ditt minne skall förbliva\\
och än åt livets bistra strid,\\
ett rosigt skimmer giva.\\
Snart tystnar allt vårt yra skämt,\\
vår sång blir stum, vårt glam förstämt.\\
O, jerum, jerum, jerum.\\
O, quae mutatio rerum!\par
\vspace{10pt}
Var äro de som kunde allt,\\
blott ej sin ära svika.\\
Som voro män av äkta halt\\
och världens herrar lika?\\
De drogo bort från vin och sång\\
till vardagslivets tråk och tvång.\\
O, jerum, jerum, jerum.\\
O, quae mutatio rerum!\\
\begin{leftborder}
\begin{tabular}{l l}
  (Filosofer)  & Den ene vetenskap och vett\\		
               & in i scholares mänger,\\
  (Jurister)   & Den andre i sitt anlets svett\\
               & på paragrafer vränger,\\
  (Teologer)   & en plåstrar själen, som är skral,\\
  (Medicinare) & en lappar hop dess trasiga fodral;\\
  (Alla)       & O, jerum, jerum, jerum,\\
               & O, quae mutatio rerum!
\end{tabular}
\end{leftborder}
\newpage
Men hjärtat i en sann student,\\
kan ingen tid förfrysa.\\
Den glädjeeld, som där har tänt,\\
hans hela liv skall lysa.\\
Det gamla skalet brustit har\\
men \textit{kärnan} finnes frisk dock kvar,\\
och vad han än må mista,\\
den skall dock aldrig brista!\par
\vspace{10pt}
Så sluten, bröder, fast vår krets,\\
till glädjens värn och ära!\\
Trots allt vi tryggt och väl tillfreds,\\
vår vänskap trohet svära.\\
Lyft bägar'n högt, och klinga vän!\\
De gamla gudar leva än\\
\revrpt bland skålar och pokaler\rpt\par
\vspace{10pt}
{\footnotesize\textit{Text: August Lindh, 1925\\ Musik:
Eugen Höfling, 1825}}\par
\vspace{10pt}
{\footnotesize\textit{Vid ``kärnan' dunkas näven i bordet en
gång. Den sista versen sjungs stående samt stående på stolen, efter
sista versen sjungits sätter sig ingen åter till bords.}}

\chapterpage{Datavetarvisor}{Datavetarvisor}
\newpage
\thispagestyle{empty}
\null
\vfill
{\noindent{\textit{\centering Fader Dator, som är i Salen,\\ helgad
    vare Din skärm,\\ tillkomme Ditt tangentbord,\\ ske Din vilja,
    så som i minnet,\\ så ock på nätet.\\ Vår dagliga utskrift giv
    oss idag\\ och förlåt oss våra misstag\\ trots att vi ej
    förlåta\\ dem som skyldiga äro.\\ Låt oss icke behöva
    vänta\\ och fräls oss ifrån dumpar.\\ Ty Salen är din\\ och
    Makten hos Handledaren\\ i Evighet.\\ UNIX.\\}}}
\vfill
\null
\newpage
\fvisa{DATAVETARNAS HÄRJARVISA}{Nu ska vi kompilera}
      {\footnotesize\textit{Melodi: Gärdebylåten}}\\
\\
Nu ska vi kompilera,\\
supa och penetrera\\
söka rekusioner ända in\\
till fixpunktssemantik.\\
Av induktion och basfall\\
blir man med lätthet asknall.\\
Hashtabeller parsas rekursivt\\
ifrån hår till häl\\
\\
{\footnotesize\textit{Text: Pär Mattsson och Magnus Ingelbo från DVL Uppsala.}}

\vspace{15pt}
\fvisa{JUBILEUMSVISA (MADHOUSE)}{En datavetare, han bor vid FooBar}\par
{\footnotesize\textit{Melodi: En sockerbagare}}\par
\vspace{10pt}
{En datavetare, han bor vid FooBar\\
sitter och kodar, vill aldrig blir klar\\
har overallen, han har långt hår.\\
har redan suttit i många år\\
Och denna dagen den skall nu firas\\
hopp in i duschen, sig själv omsvidas\\
dags att klippa sitt långa hår\\
tänk att det redan gått tretti år\par
\vspace{10pt}
{\footnotesize\textit{Text: Erika Elmquist, Kristiina Ausmees och Jens
Rosén från DV Uppsala}}\par
{\footnotesize\textit{Skriven inför DV
Uppsalas Jubileumsgasque (30 år jubiléet)}}\\
{\footnotesize\textit{Sjungs med fördel med
amerikansk brytning}}

\newpage
\svisa{THE BASIC SONG@\texttt{THE BASIC SONG}}{10 LET oss nu fatta i våra glas@\texttt{10 LET oss nu fatta i våra glas}}{\texttt{THE BASIC SONG}}
{\footnotesize\textit{Melodi: Mors lilla Olle}}\\
\\
\texttt{10 LET oss nu fatta i våra glas}\\
\texttt{20 INPUT en klunk utav det som där has}\\
\texttt{30 IF du fått nog THEN till 50 min vän}\\
\texttt{40 ELSE GOTO-baka till 10 igen}\\
\texttt{50 END}\\
\\
{\footnotesize\textit{Inledande radnummer sjungs ej.\\ I övrigt följs
    kommandon, ex. drick vid ``\texttt{INPUT en klunk}...''}}

\vspace{15pt}
\svisa{GE MIG MERA"!}{Stackars kursexaminatör}{GE MIG MERA!}
{\footnotesize\textit{Melodi: Fula gubbar}}\par
\vspace{10pt}
Stackars kursexaminatör,\\
som skall få smaka på allt som jag gör.\\
Det blir Pascal, Prolog och Lisp och C.\\
En inlupp ska du se!\\
Jag kallas hacker, freak och nörd\\
och jag kallas även störd.\\
För mig så dansar UNIX balett.\\
Jag märker ej att jag luktar svett.\\
Med en turbo-hackad fönstermiljö\\
som en manlig laddad mö.\par
\vspace{10pt}
Ge mig mera! Kompilera!\\
Runt på Polacksbacken springer jag.\\
Snart kan jag avancera.\\
Datavetare jag blir en dag.\par
\newpage
På HTML är jag autodidakt.\\
Använder News för att få kontakt.\\
Jag låter GIF-bilder bli en ångventil,\\
för lust så pubertil.\\
Och när jag en gång kommer ut,\\
blir min ingångslön till slut,\\
så hög att till och med jag får en chans.\\
Amor vill ge mig en segerkrans,\\
när jag får med mig upp på mitt rum,\\
en ny Pentium.\par
\vspace{10pt}
Ge mig mera...\par
\vspace{10pt}
{\footnotesize\textit{Text: Måns af Klercker\\Framförd på DVP
      Uppsalas Jubileumsgask (15 år jubiléet)}}

\newpage
\fvisa{DATAVISA}{Jag är helröd och helt OK}
{\footnotesize\textit{Melodi: The lumberjack song}}\par
\vspace{10pt}
Jag är helröd och är helt OK.\\
Jag jobbar hårt och jag roar mig.\par
\vspace{10pt}
Han är helröd och är helt OK.\\
Han jobbar hårt och han roar sig.\par
\vspace{10pt}
Data är ball,\\
jag kan Pascal.\\
Till FooRum vill jag gå.\\
Där träffas alla vänner\\
som är från Uppsala.\par
\vspace{10pt}
Data är ball,\\
han kan Pascal.\\
Till FooRum vill han gå.\\
Där träffas alla vänner\\
som är från Uppsala.\par
\vspace{10pt}
För han är helröd och är helt OK.\\
Han jobbar hårt och han roar sig.\par
\vspace{10pt}
Min mattebok\\
den gör mig klok.\\
Jag läser semantik.\\
Jag går på föreläsning\\
och älskar juridik.\\
\newpage
Hans mattebok\\
den gör han klok.\\
Han läser semantik\\
Han går på föreläsning\\
och älskar (\textit{förvånat}) juridik?\par
\vspace{10pt}
Men han är helröd och är helt OK.\\
Han jobbar hårt och han roar sig.\par
\vspace{10pt}
Som ekonom\\
jag blir fantom.\\
Konkurser gör mig säll.\\
Till flickor blankt jag nekar,\\
jag älskar en tabell.\par
\vspace{10pt}
(\textit{förvånad}) Som ekonom\\
han blir fantom (\textit{ursinnigt}) näääh, buuuuh!\par
\vspace{10pt}
(\textit{hurtigt}) \revrpt Men han är helröd och är helt OK.\\
Han jobbar hårt och han roar sig.\rpt\par
\vspace{10pt}
{\footnotesize\textit{Modifierad av NeNNe, dvl89, inför
    DVL Umeås 5 - års jubileum.\\ Orginalet är en
    teknologvisa från Lund.\\ (Ljusblå, KS och Umeå) har bytas mot
    (helröd, FooRum och Uppsala) för att översättas till uppländska.}}

\newpage
\fvisa{WRITE IN C}{When I find my code in tons of trouble}
{\footnotesize\textit{Melodi: Let it be}}\par
\vspace{10pt}
When I find my code in tons of trouble,\\
friends and colleagues come to me,\\
speaking words of wisdom:\\
``Write in C.''\par
\vspace{10pt}
As the deadline fast approaches,\\
and bugs are all that I can see,\\
somewhere, someone whispers:\\
``Write in C.''\par
\vspace{10pt}
Write in C, write in C,\\
write in C, oh, write in C.\\
LOGO's dead and buried,\\
write in C.\par
\vspace{10pt}
I used to write a lot of FORTRAN,\\
for science it worked flawlessly.\\
try using it for graphics!\\
Write in C.\par
\vspace{10pt}
If you've just spent nearly 30 hours\\
debugging some assembly,\\
soon you will be glad to\\
write in C.\par
\newpage
Write in C, write in C,\\
write in C, yeah, write in C.\\
Only wimps use BASIC.\\
write in C.\par
\vspace{10pt}
Write in C, write in C \\
write in C, oh, write in C.\\
Pascal won't quite cut it.\\
write in C.\par
\vspace{10pt}
Write in C, write in C,\\
write in C, yeah, write in C.\\
Don't even mention COBOL.\\
Write in C.\par
\vspace{10pt}
{\footnotesize\textit{Text: Kriston J. Rehberg}}

\newpage
\fvisa{OM EMACS}{Emacs är en stor koloss}
{\footnotesize\textit{Melodi: Rullan går}}\\
\\
Emacs är en stor koloss.\\
Tugga på, tugga på.\\
Verkar aldrig komma loss.\\
Tugga på, tugga på.\\
För att flytta på ett tecken\\
går två meg i garbage-säcken.\\
Ner i skärmens undre ran’\\
står: garbage collecting, done!\\
\\
Väntar på’n till min pension.\\
Tugga på, tugga på.\\
Emacs gör en Sparcstation\\
Tugga på, tugga på.\\
snabb som en ABC-80.\\
Högprestanda blir så smått i\\
Emacs, den ser säkert till\\
ingen dödtid blir förspilld!\\
\\
Emacs skriven är i Lisp.\\
Tugga på, tugga på.\\
Inget rimmar alls på Lisp.\\
Tugga på, tugga på.\\
Om du in i kärnan trevar\\
finn en gnu som motorn vevar.\\
Den drar Emacs runt för hand!\\
Klart att det går trögt ibland!\\
\\
{\footnotesize\textit{Text: Ingemar Ragnemalm \oldstylenums{1992}}}

\newpage
\fvisa{1515}{Höstvindar vina, ursh vilken pina}
{\footnotesize\textit{Melodi: Vårvindar friska}}\\
\\
Höstvindar vina, ursh vilken pina,\\
jag tror jag stannar hemma idag.\\
Vid terminalen, inne i salen,\\
uppe i 1515.\\
Hackar på SPARC:en, skriver i C,\\
minneshantering, vad är nu det?\\
Trycker ctrl-x, trycker ctrl-s,\\
dags då att kompilera.\\
\\
gcc gnäller, lint bara skäller,\\
“Code way too ugly, won’t even try!”\\
Ändrar ett tal här, ändrar en rad där,\\
cast’ar en liten int.\\
Prova att kompilera igen,\\
titta det gick igenom ... nej men...\\
Oj! Nej! Vad händer? Gnistrar och bränner,\\
ajdå, vi fick en core-dump.\\
\\
{\footnotesize\textit{Text: Jesper Wilhelmsson, d97\\
    \\
    ... Ja, jag vet att det ska vara C-x...}}

\newpage
\fvisa{KOMMANDOTOLKEN}{Far, jag kan inte få upp min kommandotolk}
{\footnotesize\textit{Melodi: Far, jag kan inte få upp min kokosnöt}}\\
\\
Far, jag kan inte få upp min kommandotolk.\\
Alla sätt jag prövat har vart fel;\\
jag har tryckt på knappar, klickat på varje flik,\\
och jag har sagt, “kommandotolk, tack!”,\\
men skärmen är sig lik.\\
\\
Far, jag kan inte få upp min kommandotolk,\\
inte ens med länk till satellit.\\
På varje antenn, har jag satt kopplingssladdar, men,\\
av kommandotolken min syns inte ett skit!\\
\\
Av kommandotolken min syns inte ett ski-i-it!\\
Jag ser bara Emacs, Xfig och Mosaic.\\
Säg, minns du prompten, far?\\
Glöm bort den, är du rar.\\
för av kommandotolken min syns inte ett skit!\\
\\
Far, jag kan inte få upp min kommandotolk,\\
inte ens med nitroglycerin.\\
När jag hade laddat, mycket omsorgsfullt och väl,\\
bortåt jag sprang, sen sa’ det PANG,\\
och Polacksbacken brann ner...\\
\\
Far, jag kan inte få upp min kommandotolk,\\
för jag har ingen hand med dynamit.\\
Av vårt stolta hus, det återstår bara grus,\\
och av kommandotolken min syns inte ett skit!\\
\\
Av kommandotolken min syns inte ett ski-i-it!\\
Allt det andra har blitt lika indefinit (så att säga).\\
Internet står still, vad ska jag ta mig till?\\
Av kommandotolken min syns inte ett skit,\\
ett skit,\\
ett skit,\\
ett skit,\\
Av kommandotolken min syns inte ett ski-i-it!\\
\\
{\footnotesize\textit{Text: Johan Runeson, Nollespexet DVP Uppsala,
    \oldstylenums{1996}\\ \\ (Personligen tycker jag att det ska vara
    ’Fan, jag kan inte...’ /red. anm.)}}

\vspace{15pt}
\visa{MORS LILLA DATOR}
{\footnotesize\textit{Melodi: Mors lilla Olle}}\par
\vspace{10pt}
Mors lilla dator åt skogen gick,\\
mitt i programmet så sade det klick\\
svart bidde skärmen och minnet försvann,\\
den informationen kan ingen få fram.\par
\vspace{10pt}
Brummeli-brum vad brummar där?\\
Det sprakar och gnistrar, ett jordfel det är!\\
Blixtarna blå ifrån burken det slå,\\
synd att jag nu här ensammen stå.\par
\vspace{10pt}
Hyscheli-hysch vad prasslar här?\\
Fram väller pappret ur printern där!\\
Den har fått nippran av tecken så små,\\
jag tror att jag snart hemåt skall gå.

\newpage
\fvisa{ANOTHER GLITCH IN THE CALL}{We don't need no indirection}
{\footnotesize\textit{(sung to the tune of a not so recent Pink Floyd song)}}\par
\vspace{10pt}
We don't need no indirection\\
We don't need no flow control\\
No data typing or declarations\\
Hey! Did you leave the lists alone?\par
\vspace{10pt}
All in all, it's just a pure-LISP\\
function call\par
\vspace{10pt}
We don't need no side effect-ing\\
We don't need no scope control\\
No global variables for execution\\
Hey! Did you leave those args alone?\par
\vspace{10pt}
All in all, it's just a pure-LISP\\
function call\par
\vspace{10pt}
We don't need no allocation\\
We don't need no special nodes\\
No dark bit-flipping in the functions\\
Hey! Did you leave the bits alone?\par
\vspace{10pt}
All in all, it's just a pure-LISP\\
function call
\newpage
We don't need no compilation\\
We don't need no load control\\
No link edit for external bindings\\
Hey! Did you leave that source alone?\par
\vspace{10pt}
All in all, it's just a pure-LISP\\
function call\par
\vspace{10pt}
{\footnotesize\textit{Text: DVLs tidning \#11subscript:8}}

\vspace{15pt}
\fvisa{TUPPENS KLAGOVISA}{Tuppens väckarklocka ringer sex varje morgon}
{\footnotesize\textit{Melodi: Hoppsan vilken dag}}\\
\\
Tuppens väckarklocka ringer sex varje morgon\\
då måste han gå upp till ännu en jävlig dag.\\
Han sätter sig på Teknikum och får genast ångest\\
rysningar i kroppen alla världens obehag.\\
\\
Ännu en föreläsning om fysik och sånt\\
men han tänker på en svettig datasal.\\
Tänk att sitta vid en dator i 1411\\
men miniräknarn är hans enda val.\\
\\
Tuppens verklighet det är ett liv utan glädje\\
gula overaller är ju inget man vill ha.\\
Jag borde ändå gjort som min gamla mamma sa,\\
en DVP:are é va man borde va\\
\revrpt en DVP:are é va man borde va\rpt

\newpage
\fvisa{HALLANDS-SHOT}{Ingen vet det jag vet en hemlighet}
{\footnotesize\textit{Melodi: Margareta}}\par
\vspace{10pt}
Ingen vet det jag vet en hemlighet\\
vad det var i glaset mitt som jag svept\\
kan inte hjälpa att jag känner det så här\par
\vspace{10pt}
Och den kick som jag fick gav mig allt\\
världen börjar snurra runt vad jätteballt\\
tänk att allting känns så härligt när\par
\vspace{10pt}
man dricker Rhoca-Gil ska du veta \\
då blir alla öppningar täta\\
står här bak' en dörr\\
knip' som aldrig förr\\
och kan inte pissa\par
\vspace{10pt}
pulsarna de bränner så heta\\
dricker du som jag ska du veta\\
suparna du tar\\
har du evigt kvar - inom dig\par
\vspace{10pt}
Så jag går ut igen till nästa bar\\
ännu en Hallands-shot jag genast tar\\
tänk att allting känns så härligt när\par
\vspace{10pt}
man dricker Rhoca-Gil ska du veta \\
då blir alla öppningar täta\\
står här bak' en dörr\\
knip' som aldrig förr\\
och kan inte pissa

pulsarna de bränner så heta\\
dricker du som jag ska du veta\\
suparna du tar \\
har du evigt kvar - inom dig\par
\vspace{10pt}
{\footnotesize\textit{Text: Anders Ramsell, Johan Runeson, Henrik Eriksson}}

\vspace{15pt}
\visa{DET VAR FEST NER I FOORUM}
{\footnotesize\textit{Melodi: Det var dans bort i vägen}}\par
\vspace{10pt}
Det var fest ner i FooRum på tisdagsnatten.\\
Upp i trappan gick låten av sången och skratten.\\
Det var tjo! Det var tjim! Det var hej!\\
Där satt dataloger och skåla och söp\\
och skråla och skräna tills hemåt de kröp\\
för dudeli dudeli dej. HEJ!\par
\vspace{10pt}
Där var Starka å Åhus å Renat å Bäsk,\\
där var Herrgårds och Pors som ej luktar av mäsk\\
där var OP och lite Mai Tai.\\
Där var Årsta och Vinbärs och liknande klägg,\\
där var Punsch och likör utav lingon och ägg\\
för dudeli dudeli dej. SKÅL!\par
\vspace{10pt}
{\footnotesize\textit{Text: Andreas Gref, dvl90 Uppsala}}

\newpage
\fvisa{DET SKA VA' HEMBRÄNT}{Jag har festat mycket, det har inte varit lätt}
{\footnotesize\textit{Melodi: Husvagn}}\par
\vspace{10pt}
Jag har festat mycket, det har inte varit lätt.\\
Blanda vin och whisky, det kan inte vara rätt.\\
Jag har festat på de allra konstigaste sätt,\\
men äntligen jag funnit hur man ska partajja rätt:\par
\vspace{10pt}
Det ska va hembränt, och supa till i några dar.\\
Det ska va hembränt, drick upp det sista ni har kvar.\\
Det ska va hembränt, och svina runt i overall.\\
Det ska va hembränt, drick nu vad ni tål!\par
\vspace{10pt}
I flera år så var vi lite fina där i kanten.\\
Vi smuttade på vinerna men spara aldrig panten.\\
Men sen vi lärde känna våran vackra overall,\\
så super vi nu landet runt på hemgjord alkohol.\par
\vspace{10pt}
Det ska va hembränt ...\par
\vspace{10pt}
Fem små groggar först för att komma igång med festen.\\
Fem små groggar till för att supa ikapp resten.\\
Fem små groggar sen för att kalla det en kväll.\\
Fem små groggar till så sitter vi i fyllecell.\par
\vspace{10pt}
Det ska va hembränt ...\par
\vspace{10pt}
{\footnotesize\textit{Text: Jesper Wilhelmsson, d97}}

\newpage
\fvisa{DATAVETARNAS PARADMARCH}{Ja, här kommer datavetare}
{\footnotesize\textit{Melodi: Fritjof Anderssons paradmarch}}\par
\vspace{10pt}
Ja, här kommer datavetare\\
med märken på vår dräkt.\\
Vi knappar kod, \\
vi hackar LISP\\
och surfar runt på nätet.\par
\vspace{10pt}
Vi lämnar pöbeln stående\\
och gapande förskräckt.\\
Det vi inte kan, \\
det kan ingen ann -\\
nu går vi upp i labbet.\par
\vspace{10pt}
Och när vi loggat ut,\\
ja då går vi på någon fest\\
och tar oss några bira,\\
för det gillar vi ju bäst!\par
\vspace{10pt}
För där är där vi inte är\\
och här är där vi är!\\
Vi knappar kod,\\
vi hackar LISP\\
och festar runt Sverige.\par
\vspace{10pt}
{\footnotesize\textit{Text: Selander, C93 \& Daniel ``Holm''
    Holmgren, C96.\\ Skrevs i baksätet på en Toyota Camry, på väg
    till linjekonferensen -97 i Uppsala. Omdöpt av Jubileumsgruppen
    inför 30 år jubiléet i Uppsala}}

\newpage
\fvisa{IT IS FREE}{When I find myself in front of Windows}
{\footnotesize\textit{Melodi: Let It Be}}\par
\vspace{10pt}
When I find myself in front of Windows,\\
Father Torvalds comes to me,\\
Speaking words of Linux, ``It is free!''.\\
And when hackers crack my system\\
through some strange NT obscurity,\\
he gives me the source-code, ``It is free!''.\par
\vspace{10pt}
It is free, it is free, it is free, it is free.\\
Linux is the answer, and it's free.\par
\vspace{10pt}
When the broken system crashes\\
and blue screens of death I see.\\
There is one true rescue, and it's free.\\
For though there may be doubts first\\
it will sure bring life to your PC\\
`Cause Linux is the answer, and it's free.\par
\vspace{10pt}
It is free, it is free, it is free, it is free,\\
You even get the source-code, it is free.\par
\vspace{10pt}
In the darkest night, when hacking,\\
Xfree stares right back at me,\\
Now I've installed Linux, I feel free.\\
No longer mad with my computer,\\
No more does it crash on me\\
Now I've installed Linux, it is free.
\newpage
It is free, it is free, it is free, it is free,\\
Open-sourced true power, it is free.\par
\vspace{10pt}
{\footnotesize\textit{Text: David Weinehall}}

\vspace{15pt}
\fvisa{OH, LJUVA ÖL}{Jag har nu fått en och annan öl}
{\footnotesize\textit{Melodi: Kung Louis sång}}\\
\\
Jag har nu fått en och annan öl,\\
som jag genast druckit har.\\
Ja, jag har druckit mycket jag.\\
Jag ingen droppe spar.\\
För nu vill jag bara festa,\\
tills jag far ikull.\\
Jag vill ej längre nykter va.\\
Jag vill bara vara full.\\
\\
Oh, ljuva öl, vad jag är glad åt öl.\\
För jag vill dricka den, och en till, öl.\\
Det vill jag nu, ett djur som jag,\\
det lär sig bra bli en fyllehund.\\
\\
Försök inte hindra mig gosse,\\
för nu så är det fest.\\
Och då man drycker dricka ska,\\
som gjorda är på jäst.\\
Så ge mig nu inget annat,\\
att hälla i mitt glas.\\
för nu så ska jag dricka ur,\\
och bli full som ett as.

\chapterpage{Snapsvisor}{Snapsvisor}
\section{SUPARNAS ORDNING}\index{SUPARNAS ORDNING}
\vspace{10pt}
\begin{enumerate}[leftmargin=0.65cm,topsep=0pt,itemsep=0pt,partopsep=0pt,parsep=3pt]
\item Helan
\item Halvan
\item Tersen.
\item Kvarten
\item Kvinten
\item Sexten
\item Septen
\item Rivan
\item Rafflan.
\item Rännan
\item Smuttan
\item Smuttans unge
\item Femton droppar
\item Lilla Manasse
\item Lilla Manasses Broder
\item Klämtaren
\item Kreaturens återuppståndelse
\item Den bleka dödens dryck
\item Krokofanternas intåg
\item Ett evigt liv
\end{enumerate}
\newpage
\section{ANVISNINGAR FÖR SNAPSENS TAGANDE}\index{ANVISNINGAR FÖR SNAPSENS TAGANDE}
\vspace{10pt}
\begin{itemize}[leftmargin=0cm,topsep=0pt,itemsep=2pt,partopsep=0pt,parsep=3pt,label={}]
\item Snapsen tages till fisk.
\item All mat är fisk, förutom ölsupa och pannkaka.
\item Den som äter ölsupa eller pannkaka, sup på grannens fisk.
\item Skulle grannen förtära ölsupa eller pannkaka, må även dessa rätter betraktas som fisk.
\item Desserten ätes med gaffel, så länge det går, därefter med sked.
\item Man må icke dricka ur sköljkopparna.
\item Ty den goda törsten bör släckas med ädlare fludium än tvättvatten.
\item Vatten betraktas aldrig som fisk.
\end{itemize}
\chapterpage{Ölvisor}{Ölvisor}

\chapterpage{Vinvisor}{Vinvisor}
\fvisa{GLÖGGVISA}{Vi ser det snöar}
{\footnotesize\textit{Melodi: Vi ser det snöar}}\par
\vspace{10pt}
Vi ser det snöar,\\
vi ser det snöar,\\
och glöggen vi tar.\\
Vi ser det snöar,\\
vi ser det snöar,\\
varmt i magen hurra!\\
Nu tar vi russinen fram\\
att krydda den med.\\
Sedan magen den får,\\
Och hej vad den mår.

\vspace{15pt}
\fvisa{GLÖGGVISAN}{Ack, glöggen den står på bordet}
{\footnotesize\textit{Melodi: Flickan hon går i ringen}}\\
\\
Ack, glöggen den står på bordet\\
och ångar sig varm.\\
Vi häller den ner i halsen \\
och värmer vår tarm.\\
En här, en där, utan minsta besvär,\\
den drycken är vår strupe så innerligt kär!

\newpage
\visa{SUDDA SUDDA}
{\footnotesize\textit{Melodi: Sudda sudda}}\\
\\
Sudda sudda sudda sudda bort din sura min\\
Med fyra jättestora bamseklunkar ädelt vin.\\
Munnen den ska sjunga och va' glad'.\\
För att den ska bli som den ska va'.\\
Vad häller du då bak det dolda flinet?\\
Vinet, som suddar suddar bort din sura min.

\vspace{15pt}
\fvisa{UNDULATEN}{Jag är en liten undulat}
{\footnotesize\textit{Melodi: Med en enkel tulpan}}\\
\\
Jag är en liten undulat\\
som får så dåligt med mat\\
för dom jag bor hos\\
ja dom jag bor hos\\
dom är så snåla.\\
Jag får ju fisk varenda dag\\
men det vill jag inte ha\\
jag vill ha rödvin\\
jag vill ha rödvin och gorgonzola

\newpage
\fvisa{SOLSKEN I BLICK}{Sjungom så glada till glasens klang}\\
{\footnotesize\textit{Melodi: Mors lilla Olle}}\\
\\
Sjungom så glada till glasens klang\\
vinet det har både heder och rang\\
Alla så får vi på ett ögonblick\\
rosor på kinden och solsken i blick.\\
\\
Gudarnas gåva. O, vilken fröjd.\\
Dofterna stiga mot himlens höjd.\\
Noak planterade vin och han fick\\
rosor på kinden och solsken i blick.\\
\\
Trivsamt och festligt det ska det va'\\
vinet gör människan vänlig och glad\\
vinet ger färg åt den blekaste prick\\
rosor på kinden och solsken i blick.\\
\\
Bröder och systrar nu höj ditt glas.\\
Vinet förljuvar vårt glada kalas.\\
Drick ur ditt glas och du får på ett kick\\
rosor på kinden och solsken i blick.\\
\\
Låtom oss glädjas och kasta loss.\\
Festföremålet skall säga om oss:\\
Gästerna hade, då hemåt de gick\\
rosor på kinden och solsken i blick.

\newpage
\fvisa{FRANSK VINVISA}{Feta fransyskor som svettas om fötterna}
{\footnotesize\textit{Melodi: Schuberts militärmarsch (Tomtarnas vaktparad.)}}\par
\vspace{10pt}
Feta fransyskor som svettas om fötterna,\\
de trampa druvor som sedan ska jäsa till vin.\\
Transpirationen viktig é,\\
ty den ger fin bouquet.\\
Vårtor och svampar följer mé\\
men vad gör väl dé?\par
\vspace{10pt}
För vi vill ha vin, vill ha vin, vill ha mera vin\\
även om följderna bli att vi må lida pin\par
\begin{tabular}{@{}m{0.05\linewidth}p{0.7\linewidth}@{}}
  \scalebox{1.5}{\Female} & Flaskan och glaset gått i sin\\
  \scalebox{1.5}{\Male} & Hit med vin, mera vin\\
  \scalebox{1.5}{\Female} & Tror ni att vi är fyllesvin\\
\end{tabular}\par
(\textit{skrik}) Ja, fast större!\par
\vspace{10pt}
{\footnotesize\textit{Text: Helène Derand\\Sångarstriden 1985 av
    K-LTH}}

\vspace{15pt}
\fvisa{VINETS LOV}{Läppar friska, pocka viska. Skål på Er.}
{\footnotesize\textit{Melodi: Punschen kommer}}\\
\\
Läppar friska, pocka viska. Skål på Er.\\
Strupar klara, sjunga svara, skål på Er.\\
Skål för glada minnen, skål för evig vår,\\
Inga sorger finnas mer, när vin vi får.

\chapterpage{Avecvisor}{Avecvisor}

\chapterpage{Gasquevisor}{Gasquevisor}
\visa{VÄLKOMMEN HIT}
{\footnotesize\textit{Melodi: Jänta å ja}}\par
\vspace{10pt}
Välkommen hit, tag er en bit\\
njut utav maten vi tagit hit.\\
Skoj ska vi ha, dansa vi ska\\
trivas så där förtroligt.\\
Humöret på alla toppen ska nå\\
skratta och sjunga det ska vi få.\\
Nöjda med festen hemåt vi gå.\\
I kväll ska vi ha det roligt.

\vspace{15pt}
\fvisa{TILL DRYCKERNAS ÄRA}{Vin är gott men öl är ljuvligt}
{\footnotesize\textit{Melodi: Här är gudagott att vara}}\par
\vspace{10pt}
Vin är gott men öl är ljuvligt\\
tänk vad snapsen värmer gott.\\
Kvällen våt och morgonen gruvlig.\\
Tänk att livet ändå är kort.\\
Punschen guldgul, konjak dito.\\
Vatten smakar inte gott.\\
Lär din läxa, drick ej lite.\\
Drick ur varje glas du fått.\par
\vspace{10pt}
{\footnotesize\textit{Text: Ove Andersson\\ \\ Ur spexet Jämmerdal,
    1994, Norrlands Nation i Uppsala}}

\newpage
\fvisa{VÅR UNGDOMS FAGRASTE VÅR}{Det var i vår ungdoms fagraste vår}
\vspace{10pt}
Det var i vår ungdoms fagraste vår,\\
vi drack varandra till,\\
och vi sade gutår!\\
Alla så dricka vi N.N till.\par
\vspace{10pt}
{\textbf N.N han säger inte nej därtill.}\par
\vspace{10pt}
För det var i vår ungdoms fagraste vår,\\
vi drack varandra till,\\
och vi sade gutår! 

\vspace{15pt}
\fvisa{SIFFERVISAN}{1 2 sjuttiofem 6 7}
{\footnotesize\textit{Melodi: Ritsch ratsch...}}\par
\vspace{10pt}
1 2 sjuttiofem 6 7, sjuttiofem 6 7, sjuttiofem 6 7\\
1 2 sjuttiofem 6 7, sjuttiofem 6 73\\
107 103 102\\
107 6 19 27\\
17 18 16 15\\
13 19 14 17\\
19 16 15 11\\
8 42

\newpage
\visa{JAG HAR ALDRIG VART PÅ SNUSEN}
{\footnotesize\textit{Melodi: O hur saligt att få vandra}}\par
\vspace{10pt}
Jag har aldrig vart på snusen,\\
aldrig rökat en cigarr, halleluja!\\
Mina dygder äro tusen,\\
inga syndiga laster jag har.\\
Jag har aldrig sett nåt naket,\\
inte ens ett litet nyfött barn!\\
Mina blickar går mot taket,\\
därmed undgår jag frestarens garn, halleluja!\par
\vspace{10pt}
Halleluja...\par
\vspace{10pt}
Bacchus spelar på gitarren,\\
satan spelar på sitt handklaver!\\
Alla djävlar dansar tango,\\
säg vad kan man väl önska sig mer?\\
Jo, att alla bäckar vore brännvin,\\
Fyrisån var full av bayerskt öl.\\
Konjak i varenda rännsten\\
och punsch i varendaste pöl, mera öl!\par
\vspace{10pt}
Mera öl...\par
\vspace{10pt}

\newpage
\fvisa{KAMELEN PÅ KANELEN}{Om jag vore ett tvåpuckligt djur}
{\footnotesize\textit{Melodi: Old McDonald}}\par
\vspace{10pt}
Om jag vore ett tvåpuckligt djur\\
Eta-eta-nol\\
Kunde jag super i ur och skur\\
Eta-eta nol\\
Fett, fett, fett, bärverfett\\
Bäverfett gör livet lätt\par
\vspace{10pt}
När man går i Kemikums öken\\
Eta-eta-nol\\
Vill man gärna ta till kröken\\
Eta-eta-nol\\
Fett, fett, fett, ...\par
\vspace{10pt}
När jag spriten lagrat upp\\
Eta-eta-nol\\
Klår jag lättans upp en mupp\\
Eta-eta-nol\\
Fett, fett, fett, ...\par
\vspace{10pt}
Nu när nubben slinker ner\\
Eta-eta-nol\\
Jag redan längar efter mer\\
Eta-eta-nol\\
Fett, fett, fett, ...

\newpage
\visa{DOM SOM ÄR NYKTRA}
{\footnotesize\textit{Melodi: Du är den ende}}\par
\vspace{10pt}
Dom som är nyktra har inte så roligt\\
dom har bara ansvar och inte nåt tjolit-\\
an-lej-faderulla, men vi som är fulla\\
vi har bara kul nästan jämt.\par
\vspace{10pt}
Det sägs att en männ'ska kan va' utan brännvin\\
det stämmer kan hända men se blott den min\\
som pryder en absolutist den är djävligt trist\\
därför sjunger vi så:\par
\vspace{10pt}
Dom som är...

\vspace{15pt}
\fvisa{SÅ LÄNGE}{Så länge rösten är mild}
{\footnotesize\textit{Melodi: Så länge skutan kan gå}}\\
\\
Så länge rösten är mild,\\
så länge ingen är vild,\\
så länge spegeln på väggen ger halvskaplig bild.\\
Så länge alla kan stå,\\
så länge alla kan gå,\\
så länge alla kan tralla så fyller vi på.\\
För vem har sagt att just du kom med storken,\\
för att bli glad av att lukta på korken?\\
Men kring vårt bord här nånstans\\
vi höjer bägaren med glans,\\
och låter vinet gå ner med en yrande dans.

\newpage
\visa{VI SOM ÄR NYKTRA}
{\footnotesize\textit{Melodi: Du är den ende}}\par
\vspace{10pt}
Vi som är nyktra vi har faktiskt roligt.\\
Jo visst har vi ansvar, men minst lika tjolitt-\\
anlej faderulla som ni som är fulla\\
som tror ni har kul nästan jämt.\par
\vspace{10pt}
Men tänk då efter uppå dagar efter\\
de dagar med fester, med smärtsamma rester\\
utav eran hjärna, nog tycker ni gärna:\\
``Att va' nykterist är nå't visst.''\par
\vspace{10pt}
Vi som är nyktra, vi har bara roligt.\\
Imorr'n kan vi återigen ha det tjolitt-\\
anlej faderulla, men ni som var fulla\\
aj, aj, aj, det är väl för trist.

\vspace{15pt}
\visa{TAG NU DJUPA KLUNKAR}
{\footnotesize\textit{Melodi: Lili Marlene}}\par
\vspace{10pt}
Tag nu djupa klunkar,\\
sjung och drick, var glad.\\
Innan gänget lunkar\\
på festprisseparad.\\
Vad gör väl det om våra ben\\
då vingla lätt på gatans sten.\\
När vi beger oss hem.\\
Vi undrar blott till vem.

\newpage
\visa{JAG VAR FULL EN GÅNG}
{\footnotesize\textit{Melodi: Flottarkärlek}}\par
\vspace{10pt}
Jag var full en gång för länge sen,\\
på knäna kröp jag hem.\par
\vspace{10pt}
Varje dike var för mig ett vilohem.\\
I varje skåp, i varje vrå\\
hade jag en liten vän,\\
ifrån renat upp till 96 procent, procent.\par
\vspace{10pt}
Jag var full...\\
Och som sällskap hade jag en elefant.\\
Elefanten spruta vatten,\\
men jag trodde det var vin,\\
sedan dess har alla kallat mig för svin, mera vin.\par
\vspace{10pt}
Jag var full...\\
Och som sällskap hade jag en elefant.\\
Elefanten spruta vatten,\\
men jag trodde det var öl,\\
sedan dess har alla kallat mig för knöl, mera öl.\par
\vspace{10pt}
Jag var full...\\
Och som sällskap hade jag en elefant.\\
Elefanten spruta vatten,\\
men jag trodde det var sprit\\
sedan dess har alla kallat mig för skit, mera sprit

\newpage
\fvisa{LAMBO}{Sätt nu glaset till din mun"!}
\vspace{10pt}
Sätt nu glaset till din mun!\\
Tjofaderitan, lambo!\\
Och drick ur din fylle-, fyllehund!\\
Tjofaderitan, lambo!\\
Se hur dropparna i glaset\\
rinner ner i fylleaset.\\
Lambo hej, lambo hej, \\
tjofaderitan lambo!\\
\\
\textbf{Jag nu glaset druckit har!}\\
Tjofaderitan lambo\\
\textbf{Ej en droppe finnes kvar!}\\
Tjofaderitan lambo.\\
\textbf{Som bevis jag nu vill vända\\
glaset på dess ``rätta'' ända.}\\
Lambo hej, lambo hej\\
tjofaderitan lambo.\\
\\
Ja, han kunde konsten,\\
han var en äkta fylle-, fyllehund.\\
Så låt oss gå till nästa man\\
och se vad han förmår\\
\\
(\textit{En liten variant})\\
\textbf{Som bevis jag nu vill föra\\
glaset till min grannes öra}

\newpage
\fvisa{NIKOLAJEV}{Jag heter Nikolajev}
{\footnotesize\textit{Melodi: Sovjetunionens nationalsång}}\par
\vspace{10pt}
Jag heter Nikolajev\\
och kommer från Sovjet.\\
Jag flyger runt jorden\\
i min rymdraket.\\
Och där ska jag stanna i 84 varv\\
för det har Chrustjev sagt\\
men det tycker jag är larv.\par
\vspace{10pt}
Jag längtar hem, hem till min planet.\\
Till fru och barn därhemma i Sovjet.\\
Men mest utav allt längtar jag till ett rum\\
med ett hjärta på dörren.\\
Jag längtar hem till min planet\\
till fru och barn därhemma i Sovjet\par
\vspace{10pt}
Min kapsel innehåller\\
många instrument.\\
Ja, mycket av sådant\\
som ännu ej är känt.\\
Men lika förbannat, vad du än tror:\\
Jag glömde gå på muggen\\
innan jag for.\par
\vspace{10pt}
Jag längtar hem...
\par
\vspace{10pt}
{\footnotesize\textit{Uppsala teknologernas signarturmelodi, datavetare sitter ner när den sjungs.}}

\newpage
\fvisa{MINNE}{Minne"! Jag har tappat mitt minne}
{\footnotesize\textit{Melodi: Memory ur musikalen Cats}}\\
\\
Minne! Jag har tappat mitt minne\\
är jag svensk eller finne?\\
Kommer inte ihåg\\
\\
Inne! Är jag ute eller inne?\\
Jag har luckor i minne\\
såna där små alkohol.\\
\\
Men besinne, man tätar det med brännvin \\
som man får fastän minnet \\
och helan den går.\\
\\
{\footnotesize\textit{Text: Bosse Österberg}}

\vspace{15pt}
\visa{DET VAR LÄNGE SEN}
{\footnotesize\textit{Melodi: Det var länge sin jag plocka några blommor}}\par
\vspace{10pt}
Det var länge sen jag plocka' några blommor.\\
Det var länge sen jag tog några poäng.\\
Det var länge sen jag handla' på systemet.\\
Det var länge sen jag fick en tjej/kille i säng.\par
\vspace{10pt}
Men å andra sidan bränner jag ju hemma,\\
och klarar kärleken alldeles för mig själv.\\
Det var länge sen jag plocka' några tentor\\
men å andra sidan går de om igen.

\newpage
\fvisa{KALMAREVISAN}{För uti Kalmare stad}
\vspace{10pt}
\textbf{För uti Kalmare stad}\\
ja där finns det ingen kvast\\
förrän lördagen.\\
\textbf{Hej dick}\\
Hej dack\\
\textbf{Jag slog i}\\
och vi drack\\
\textbf{Hej dickom dickom dack}\\
hej dickom dickom dack.\\
För uti Kalmare stad\\
ja där finns det ingen kvast\\
förrän lördagen.\\
\\
\revrpt \textbf{När som bonden kommer hem}\\
kommer bondekvinnan ut\rpt\\
och är stor i sin trut\\
\textbf{Hej dick...}\\
\\
\revrpt \textbf{Var är pengarna du fått?}\\
Jo, dom har jag supit opp!\rpt\\
Uppå Kalmare slott.\\
Hej dick...\\
\\
\revrpt \textbf{Jag skall mäla dig an}\\
för vår kronbefallningsman\rpt\\
Och du skall få skam\\
\textbf{Hej dick...}
\newpage
\revrpt \textbf{Kronbefallningsmannen vår}\\
satt på krogen i går\rpt\\
Och var full som ett får.\\
\textbf{Hej dick...}\\
\\
{\footnotesize\textit{Text: G S Kallstenius}}\\
\\
(\textit{Tre små tillägg})\\
\revrpt \textbf{Va' sa' bonnen ha te' mat?}\\
Sura sillar och potat\rpt\\
det blir sillsallat.\\
\textbf{Hej dick...}\\
\\
\revrpt \textbf{Säg var är din lab-rapport}\\
Jo, den har jag supit bort!\rpt\\
För den var så kort!\\
\textbf{Hej dick...}\\
\\
\revrpt\textbf{Ordföranden vår}\\
var på 100-kör igår.\rpt\\
Dom bar ut han på bår!\\
\textbf{Hej dick...}\\
\\
{\footnotesize\textit{I Linköping är det inget ``Hej dick...''-ande på
    första versen, dock reprisande som följande verser.}}

\newpage
\fvisa{VIKINGEN}{En viking älskar livets vann}
{\footnotesize\textit{Melodi: When Johnny comes marching home}}\par
\vspace{10pt}
En viking älskar livets vann\\
Hurra, hurra!\\
Den hastigt i mitt svalg försvann\\
Hurra, hurra!\\
Till kalv, till oxe, till fisk, till fläsk,\\
när kärringar bara dricker läsk.\\
Ja, då vill alla vikingar ha en bäsk.\par
\vspace{10pt}
När bäsken småningom är slut\\
Tragik, tragik\\
Då bäres varje viking ut\\
som lik, sig lik\\
Och sen, om vi vaknar, vi sjunger en bit,\\
sen korkar vi upp Skånes Aquavit\\
Skål för alla vikingar som kom hit!\par
\vspace{10pt}
{\footnotesize\textit{Sångarstriden 1981 av E-LTH}}\par
\vspace{10pt}
Och datavetaren minsann,\\
vill också ha\\
en bäsk till maten då och då\\
det sitter bra\\
När datavetarens väg är klar\\
han lämnat inte en droppe har\\
skål för alla vikingar som finns kvar\par
\vspace{10pt}
{\footnotesize\textit{Text: Carl Leonardsson, Anders Steinrud, Niclas Stensbäck, Samuel Strand}}

\newpage
\visa{WE ARE SAILING}
%% {\footnotesize\textit{Melodi: Sailing}}\\
%% \\
\vspace{10pt}
We are sailing, we are sailing,\\
through the water, across the sea.\\
We are sailing, lugna waters,\\
to be ute, to be free.
\vspace{12pt}\\
\begin{tabular}{@{}m{0.1\textwidth}p{0.9\textwidth}@{}}
  \scalebox{3}{\Male} & \specialcell{I am länsing, I am glänsing,\\
    forsing framåt, mera vind!\\
    Elva meter, fulla segel,\\
    that's my life, my dear wife.}
\end{tabular}
\vspace{10pt}\\
\begin{tabular}{@{}m{0.1\textwidth}p{0.9\textwidth}@{}}
  \scalebox{3}{\Female} & \specialcell{I am freezing, it is gunging,\\
    båten laying, oh my God!\\
    Children crying, starta motor!\\
    Oh, where are we? I'll go home.}
\end{tabular}
\vspace{10pt}\\
\begin{tabular}{@{}m{0.1\textwidth}p{0.9\textwidth}@{}}
  \scalebox{3}{\Male} & \specialcell{Can't you hear me, can't you hear me?\\
    Tyst du stör mig, I am busy.\\
    I'm kappsegling, med en maxi,\\
    to be winner, or to die.}
\end{tabular}
\vspace{10pt}\\
\begin{tabular}{@{}m{0.1\textwidth}p{0.9\textwidth}@{}}
  \scalebox{3}{\Female} & \specialcell{Can't you hear me, can't you hear me?\\
    Reva segel, sun is down.\\
    Går på grundet, båt is läcking.\\
    Jag tar bussen in to town.}
\end{tabular}
\vspace{7pt}\\
We are sailing, we are sailing,\\
home again, across the sea.\\
We are flying, forever trying,\\
to be happy, to be free.

\newpage
\fvisa{HÄRJAREVISAN}{Liksom våra fäder, vikingarna i Norden}
{\footnotesize\textit{Melodi: Gärdebylåten}}\par
\vspace{10pt}
Liksom våra fäder, vikingarna i Norden.\\
Drar vi riket runt och super oss under borden.\\
Brännvinet har blitt ett elixir,\\
för kropp såväl som själ.\\
Känner du dig liten och ynklig på jorden,\\
växer du med supen och blir stor ut i orden.\\
Slå dig för ditt håriga bröst och bli \\
en man från hår till häl.\par
\vspace{10pt}
Ja, nu skall vi ut och härja,\\
supa och slåss och svärja,\\
bränna röda stugor, slå små barn och säga fula ord.\\
Med blod skall vi stäppen färga.\\
Nu äntligen lär ja'\\
kunna dra någon riktig nytta\\
av min Hermodskurs i mord.\par
\vspace{10pt}
Hurra, nu ska man äntligen få röra benen,\\
hela stammen jublar och det spritter i grenen.\\
Tänk att än en gång få spränga fram\\
på Brunte i galopp!\\
Din doft o kära Brunte är trots sin brist i hygienen,\\
för en vild mongol minst lika ljuv som syrenen.\\
Tänk att på din rygg få rida runt\\
i stan och spela topp.\par
\vspace{10pt}
Ja, nu skall vi ut och härja...\par
\vspace{10pt}
Ja, mordbränder är klämmiga ta fram fotogenen\\
och eftersläckningen tillhör just de fenomenen\\
inom brandmansyrket som jag tycker\\
det är nån nytta med.\\
Jag målar för mitt inre upp den härliga scenen;\\
Blodrött mitt i brandgult, ej ens prins Eugen en\\
lika mustig vy kan måla, \\
ens om han målade med sked.\par
\vspace{10pt}
Ja, nu skall vi ut och härja...\par
\vspace{10pt}
{\footnotesize\textit{Text: Hans Alfredsson och Levin\\
Ur Lundaspexet Djingis Kahn 1954, första versen ej ur originalet.}}\par
\vspace{10pt}
\index{DATAVETARNAS HÄRJARVISA}
Nu ska vi kompilera,\\
supa och penetrera\\
söka rekusioner ända in\\
till fixpunktssemantik.\\
Av induktion och basfall\\
blir man med lätthet asknall.\\
Hashtabeller parsas rekursivt\\
ifrån hår till häl\par
\vspace{10pt}
{\footnotesize\textit{Text: Pär Mattsson \& Magnus Ingelbo från DVL
    Uppsala.}}

\newpage
\fvisa{SPRITBOLAGET}{Till spritbolaget ränner jag}
{\footnotesize\textit{Melodi: Emil i Lönneberga}}\\
\\
Till spritbolaget ränner jag,\\
och bankar på dess port.\\
Jag vill ha nå't som bränner bra,\\
och får mig skitfull fort.\\
Expediten sade goda',\\
hur gammal kan min herre va'.\\
Har du nåt leg, ditt fula drägg,\\
kom hit igen när du har fått skägg.\\
\\
Nej, detta var ju inte bra,\\
jag skall bli full ikväll.\\
Då plötsligt en ide jag fick,\\
de har ju sprit på Shell.\\
Många flaskor stod där på ra'.\\
Så nu kan jag bli full och gla'.\\
Den röda drycken åkte ner,\\
nu kan jag inte titta mer.\\
\\
{\footnotesize\textit{Text: Göran Svensson \\ \\ Sångarstriden 1989 av
    E-LTH.}}

\newpage
\visa{SJUNG OM FRU SVENSSONS}
{\footnotesize\textit{Melodi: Studentsången}}\\
\\
Sjung om fru Svenssons lyckliga karl,\\
låt honom plöja i ungdomens fåror.\\
Fem gamla hjärtan i sprit har jag.\\
Å' en ljus elefant i ett snår.\\
Inga stoppar den,\\
i vårat linneskåp.\\
Loppor tär vår vän,\\
som idisslar en sko\\
när vi snyta en rund liten hund.\\
\revrpt Där den här lilla bagaren bor.\rpt\\
Hursa?\\
\\
{\footnotesize\textit{Text: Povel Ramel, 1948}}

\vspace{15pt}
\fvisa{VAR INTE BLYG}{Jag har inte fått nån sup}
{\footnotesize\textit{Melodi: Sov du lilla videung}}\\
\\
Jag har inte fått nån sup,\\
därför är jag blyger.\\
Uti magens trånga djup,\\
hämningarna smyger.\\
Men om blott jag får en tår,\\
hämningen ur kroppen går,\\
rodnar inte mera,\\
väntar blott på flera.

\newpage
\visa{VI KOM' FRÅN NORRLAND}
{\footnotesize\textit{Melodi: Vi bor på landet}}\par
\vspace{10pt}
Vi kom' från Norrland. Vi käkar myggen\\
som vi fångar på kalhyggen.\\
Och vi lägg' dom uppå smörgåsen\\
som vi ät' med:\\
\hspace*{25pt} laxen,\\
\hspace*{25pt} renen,\\
\hspace*{25pt} sillen,\\
\hspace*{25pt} smöre',\\
\hspace*{25pt} nubben,\\
\hspace*{25pt} kåtan,\\
\hspace*{25pt} pölsan,\\
\hspace*{25pt} pären,\\
\hspace*{25pt} fjälle',\\
\hspace*{25pt} filen,\\
\hspace*{25pt} Åkes hammock\\
 och det gör vi varje dag.\par
\vspace{10pt}
Han hette Åke och var från Tåsjö.\\
Han had' en gris som var från Blåsjö.\\
Som han lade uppå smörgåsen\\
som han käkade med laxen...\\
...och det gör han varje dag.
\newpage
Sen har vi Herman, han kom' från Böle.\\
Han hade Norrlands minsta hemman.\\
Han gick och kratte och kärringa skratte\\
så han la'na på smörgåsen\\
och han käka'na med laxen...\\
...och det gör han varje dag.\par
\vspace{10pt}
Sen har vi Göran, han kom' från Pite'\\
Har ingen hjärna men det skit' vi i.\\
För han la den uppå smörgåsen\\
som han åt med laxen...\\
...och det gör han varje dag.\par
\vspace{10pt}
Vi kom från Norrland. Ssss!\\
Vi käkar myggen Ssss!\\
som vi fångar på kalhyggen. Ssss!\\
Och vi lägg' dom Ssss!\\
uppå smörgåsen Ssss!\\
som vi ät' med laxen...\\
...och det gör vi varje daaag!

\newpage
\fvisa{OFVANDAHLS}{Hundra år sen ungefär}
{\footnotesize\textit{Melodi: Ovan där...}}\par
\vspace{10pt}
Hundra år sen ungefär\\
Ofvandahls han sa så här:\\
``Skriva dikter ger mig ingen\\
peng i pung, nej blott misär.\\
Det nåt annat måste bli,\\
varför inte bageri,\\
runda bullar små mig gör till millionär.''\par
\vspace{10pt}
Ofvandahls, kaffe med avec\\
Ofvandahls, Landings stora skräck,\\
Ofvandahls Napoleon canapé\\
ger oss krafter nog att vandra vidare.\par
\vspace{10pt}
Under åren som har gått,\\
många styrketåren fått\\
och en del av dem på Ofvandahls,\\
när det känns trist och grått.\\
Efter nattens hårda slit, \\
när det ej finns mera sprit,\\
då man tar en kaffetår och lever upp.\par
\vspace{10pt}
Ofvandahls, kaffe...

\newpage
\fvisa{VÅRSÅNG}{Glad som en viking i offerlunden}
\vspace{10pt}
Glad som en viking i offerlunden\\
Full som en pumpa på sista april\\
Punschen i magen och lever för stunden\\
Så länge som studiestödsnämnden blott vill\par
\vspace{10pt}
Ta en poäng och sen ta sig en bläcka\\
Röja och ramla på gator och torg\\
Skändligen stans innevånare väcka\\
Slumra dan efter i lärdomens borg\par
\vspace{10pt}
Se hur med fransiga frackarna på\\
gossarna små hoppa och slå\par
\vspace{10pt}
Butter och bakis mot Heimstaden hasta\\
Lever på blodpudding pilsner och pasta\\
Och höja ett glas och sen höja ett glas\\
Som ett glatt och bekymmersfritt as.

\newpage
\fvisa{INTEGRALVISAN}{En liten enkel integral}\\
{\footnotesize\textit{Melodi: Med en enkel tulipan}}\\
\\
En liten enkel integral\\
ett vektoranalystal\\
ni har besväret, ni har besväret\\
att derivera.\\
\\
Men tar man Stokes sats däruppå\\
så blir det så enkelt så\\
att integralen, att integralen\\
evaluera.\\
\\
Och rotationen, den integreras\\
sen över ytan utav en boll,\\
koordinaterna transformeras,\\
så integranden blir bara noll.\\
\\
En liten enkel integral\\
ett vektoranalystal\\
kan va så djävlig att man ej hinner\\
med något mera.

\newpage
\fvisa{LAPPSÅNGEN}{Jag är lapp och jag har lite renat}
{\footnotesize\textit{Melodi: Vid foten av fjället}}\\
\\
Jag är lapp och jag har lite renat\\
det är hemgjort, ja tacka för det,\\
ty att köpat blir dyrt så förbenat\\
för då kommer ju frakterna till.\\
\\
Här vid foten av snöklädda fjället\\
finns en plats dit jag drar varje vår\\
det är faktiskt nåt särskilt med stället\\
det är där apparaten min står.\\
\\
Men en natt var där odjur i trakten\\
där var skära elefanter och möss.\\
Och till trolltrummans ljud upptogs jakten\\
till apparaten så sa jag ajöss.\\
\\
Men när jag kom tillbaka till fjället\\
fanns av apparaten inte ett spår\\
ty se länsman han hittade stället\\
och i finkan så ensam jag går.

\newpage
\fvisa{LINGONBEN}{Bluff och Spark och Tork och Kvark}
\vspace{10pt}
Bluff och Spark och Tork och Kvark\\
voro sex små dvärgar.\\
En var ful och en var glad \\
och en var dum i huv'et.\\
Hej, sa Tork till lille Kvark,\\
känner du igelkotten Pilt?\\
Han som har varit i Paris?\\
Ja, det gjorde Ivar.\\
Hör du hans lilla runda tass\\
när som han trippar på sitt pass:\\
Tripp och trapp och trypa.\\
Se hans lilla piga!\par
\vspace{10pt}
Tomtefar i skogens brus\\
sitter som ett päron.\\
Han har inget eget hus \\
allt i sin stora näsa.\\
Söt och blöt är sagans fé.\\
Trollen är bjudna hit på te.\\
Det lilla trollet! Pass för det!\\
Nu ska mormor bada.\\
Väva och spinna natten lång.\\
Prinsen är här i fjorton språng.\\
Hopp och hipp och huppla.\\
Hästen heter Sverker!\par
\newpage
Stora slottet Drummeldimp\\
ligger bortom fjärran.\\
Dit får ingen komma in\\
som ej kan baka struvor.\\
Gyllenkrull och sockertipp.\\
Kom ska vi dansa häxan våt!\\
Vill du mig här, så har du nåt.\\
Sov du lilla tryne.\\
Kungen är full av stock och sten.\\
Skogen är full av lingonben.\\
Pär är full av tomtar.\\
Hur ska Lillan orka?\par
\vspace{10pt}
{\footnotesize\textit{Text \& Musik: Povel Ramel, 1957\\ En
    ``struva'' är ett flottyrkokt bakverk}}

\newpage
\fvisa{PROVRÖRSBARNET}{Sov du lilla provrörsgull}
{\footnotesize\textit{Melodi: Videvisan}}\\
\\
Sov du lilla provrörsgull.\\
Nu ska Lillan sova.\\
Far din han var en ampull.\\
Vetenskapens gåva.\\
Far han fanns på institut.\\
Ur ett kylskåp togs han ut.\\
Låg där på förvaring.\\
Sov nu lilla raring.\\
\\
Jag har sluppit bli förförd,\\
skämd och chikanerad.\\
Far din, han var stamboksförd,\\
silad och filtrerad.\\
Inga kriminella drag\\
ärvda enligt Mendels lag\\
fanns i genotypen.\\
Vyssja fenotypen.\\
\\
Lugn och lycklig är din mor.\\
Vyssja lilla vännen.\\
Konkurrensen blev för stor\\
utan hjälp av männen.\\
Utan gråt och utan skratt\\
fick jag dig min lilla skatt.\\
Jag har valt tekniken\\
framför erotiken.
\newpage
Sov mitt lilla frökenbarn,\\
skönt att ha varandra.\\
Skönt att slippa dras med karl'n\\
som så många andra.\\
Oanfrätt är min moral.\\
De som sprider lömskt förtal\\
tvangs att sträcka vapen.\\
Leve vetenskapen!\\
\\
{\footnotesize\textit{Ur revyn “Nachtspiel”, 1951}}

\vspace{15pt}
\fvisa{VISA FÖR OMUSIKALISKA}{Det fanns brännvin i flaskan när vi kom}\\
{\footnotesize\textit{Melodi: Ja tack, gärna det!}}\\
\\
Det fanns brännvin i flaskan när vi kom\\
(\textit{Dunk}, \textit{dunk})\\
När vi gick\\
(\textit{Dunk}, \textit{dunk})\\
Var flaskan tom\\
(\textit{Klunk}, \textit{klunk})\\
\\
Vi var klara i knoppen när vi kom\\
(\textit{Dunk}, \textit{dunk})\\
När vi gick\\
(\textit{Dunk}, \textit{dunk})\\
Var det tvärtom\\
(\textit{Klunk}, \textit{klunk})\\
\\
{\footnotesize\textit{Vid varje ``klunk'' tas en ordentlig klunk sprit (eller grogg för de svagare själarna)}}

\newpage
\fvisa{STÖRTHÄRLIGT FULL}{Nu har alla lämnat festen}
{\footnotesize\textit{Melodi: Fat Mammy Brown}}\par
\vspace{10pt}
Nu har alla lämnat festen\\
och jag sitter ensam kvar\\
bland groggar, pilsnerflaskor\\
i en sönderslagen bar.\\
Första pilsnerflaskan tog jag\\
vid min frukost klockan fem\\
och nu sitter jag och väntar\\
på att bli buren hem.\par
\vspace{10pt}
För jag är störthärligt full.\\
Jag ramlar mest omkull.\\
Jag ser skära elefanter \\
som har jättekonstig ull.\\
Ja, jag är störthärligt full\\
och ramlar mest omkull.\\
Det är präktigt, härligt;\\
supa och va' full.\par
\vspace{10pt}
Ifrån festen minns jag inget,\\
men mitt öga blev visst blått.\\
Det måste jag ha fått\\
när någon kastat en karott\\
full med vispgrädde och fimpar\\
och en okammad peruk,\\
eller också när jag stod\\
i moraklockan och var sjuk.
\newpage
För jag är...\par
\vspace{10pt}
Nästa morgon när jag vaknar\\
med en bergsborr i min kropp\\
och sandpapper på tungan.\\
Jag vill ej stiga opp.\\
Mina armar dom känns tunga\\
och min näsa den är sne'.\\
Då raglar jag till köket\\
för en återställare.\par
\vspace{10pt}
För jag är...

\vspace{15pt}
\fvisa{KANTOM STUDIOSIJ}{Kantom studiosij extrabon sjur.}\\
{\footnotesize\textit{Melodi: Studentsången}}\\
\\
Lassom galejan in spring juvenar.\\
Noch funkar kordan kum san bravur\\
Kaj futura blondina est var\\
Nolla furii in va psykosan sit\\
Esperan vmi promissan va kredit\\
Kum voj knopa bandage in plantage\\
Kvo dullkissan diploma florit\\
Kvo dullkissan diploma florit\\
HOJLA!\\
\\
{\footnotesize\textit{In transpiranto par Ludviko Hagbaldo\\ Ur
    Grönköpings veckoblad.}}

\newpage
\fvisa{RECEPTET}{19 hönor, bruna bönor}
{\footnotesize\textit{Melodi: Gubben Noak}}\par
\vspace{10pt}
19 hönor, bruna bönor\\
18 kilogram\\
17 stycken sillar ner i magen trillar\\
16 läckra, spröda, smäckra rostade små lamm\par
\vspace{10pt}
15 bullar sedan rullar ner i halsen lätt\\
14 omeletter\\
13 fläskkotletter\\
Jag bland rätter nästan sätter dem som nummer 1\par
\vspace{10pt}
12 glas dricka kan jag slicka i mig när jag vill\\
11 små potäter jag till middag äter\\
10 gäddor\\
9 skäddor\\
8 kvistar dill\par
\vspace{10pt}
7 små harar jag väl klarar, pannekakor\\
6\\
5 små gäss med krås till\\
4 torskar, sås till\\
3 fasaner\\
2 bananer och till slut\\
1 älg. HEJ!

\newpage
\fvisa{GÄLLIVAREVISAN}{Dänkte på lörda skulle fara}\\
\vspace{10pt}\\ 
Dänkte på lörda skulle fara\\
in i Gällivare på Karakkatorg.\\
Där pjuda dricka brännvinsflaskor\\
gott som sälva satan voj, voj.\\
Det var liga liven, bara slagsmål ma niven\\
åka bolisstationen, vara jävliga fasonen.\\
Ligga inne alva natten, leva limpa å vatten\\
komma ut morrokröken, säja nappast då ajöken.\\
\\
Men dängte perkele anamma\\
ganse vara samma åka Vitafors.\\
Där sänner till en gammal licka\\
som jag prukar likka förståss.\\
De prukar pli ganska sällan, fara hälsa på fjällan\\
å ganse få lite mellan, bara inte vastna fällan.\\
De är glart man riskera, ingentin reflektera\\
bara inte nå mera, då måste gå å operera.\\
\\
Men nu jak sluta erotiken\\
öppna spritfabriken, sälja akvaviten\\
å komma lansfiskalen nära,\\
fråka om jak pära spriten.\\
De vara jävliga token, pliva hånkad av snoken\\
draka satan åt skogen, åka Luleå kroken.\\
Sitta inne halva åren, pliva krå uti håren\\
se av brännvin int spåren,\\
å komma bakas först på våren.

\newpage
\fvisa{EN GLAD CALYPSO OM VÅREN}{Jag dansar runt och jag sjunger strunt}
\vspace{10pt}
Jag dansar runt och jag sjunger strunt\\
och jag är visst lite i hatten.\\
Tralalliralla i månens sken\\
där jag vandrar hemåt i natten.\\
Vart gick dom andra, var blev dom av?\\
Jag är ensam här med min flaska.\\
Tralalliralla, men det är ju kul\\
att i vattenpussarna plaska.\\
\\
Tralla-lalla-lalla-lalla.\\
Lalla-lalla-la-la.\\
Tralla-lalla-lalla-lalla-la.\\
Lalla-lalla-la-la.\\
\\
Så utmed dikena plaskar jag\\
på min stolta väg ifrån festen.\\
Tralalliralla jag trillar visst,\\
men det gör detsamma föresten.\\
Jag dansar långdans med alla trän,\\
så att mossan ryker i snåren.\\
Tralalliralla med rönn och en,\\
i en glad calypso om våren.\\
\\
Tralla...
\newpage
Jag dansa' ut på ett fält förut,\\
fullt av is som låg där och blänkte.\\
Tralalliralla precis som om\\
det var frost och snö, och jag tänkte:\\
Vad tusan frosten är svår i maj.\\
Oh, det lät som glas när jag trampa.\\
Å jädrans anamma vad jag blev skraj\\
för jag runt i drivbänkar trampa.\\
\\
Tralla...\\
\\
Men tjo, vad jag är glad ändå\\
fast jag trampat uti rabatten.\\
För det är härligt när det är vår\\
och jag dansar hemåt i natten.\\
Å, titta snart är det ljusan dag\\
alla fåglar sjunga i snåren.\\
Kom med och dansa med mig ett slag\\
i en glad calypso om våren.\\
\\
Tralla...\\
\\
{\footnotesize\textit{Text och Musik: Olle Adolphson}}


\chapterpage{Säsongsvisor}{Säsongsvisor}
\fvisa{LÄNGTAN TILL LANDET}{Vintern rasat ut bland våra fjällar}
\vspace{10pt}
Vintern rasat ut bland våra fjällar\\
drivans blommor smälta ner och dö.\\
Himlen ler i vårens ljusa kvällar,\\
solen kysser liv i skog och sjö.\\
Snart är sommarn här! I purpurvågor \\
guldbelagda azurskiftande,\\
ligga ängarne i dagens lågor,\\
och i lunden dansa källorne.\\
\\
Ja, jag kommer! Hälsen glada vindar\\
ut till landet, ut till fåglarne,\\
att jag älskar dem, till björk och lindar,\\
sjö och berg, jag vill dem återse,\\
se dem än, som i min barndoms stunder,\\
följa bäckens dans till klarnad sjö,\\
trastens sång i furuskogens lunder,\\
vattenfågelns lek kring fjärd och ö.\\
\\
{\footnotesize\textit{Text: Herman Sätherberg\\ Musik: Otto Lindblad}}

\newpage
\visa{SKÖNA MAJ}
\vspace{10pt}
Sköna maj välkommen till vår bygd igen,\\
sköna maj välkommen våra lekars vän.\\
Känslans gudaflamma väcktes vid din ljusning.\\
Jord och skyar stamma, kärlek och förtjusning,\\
sorgerna flyr för våren, glädje ler ur tåren.\par
\vspace{10pt}
Sköna maj välkommen till vår bygd igen,\\
Sköna maj välkommen våra lekars vän.\\
Känslans gudaflamma väcktes vid din ljusning.\\
Nu ur lundens sköte och ur blommande knopp,\\
stiga dig till möte, glada offerhopp.\par
\vspace{10pt}
Sköna maj välkommen till vår bygd igen,\\
Sköna maj välkommen våra lekars vän.\\
Känslans gudaflamma väcktes vid din ljusning.\\
Hör med tacksam tunga. Tusen fåglar sjunga,\\
liksom vi, liksom vi.\\
Välkommen sköna maj, välkommen sköna maj.\par
\vspace{10pt}
{\footnotesize\textit{Text: J.L. Runeberg}}

\newpage
\visa{SUMMARN KUMMARN}
\vspace{10pt}\\
Summarn kummar me sol,\\
kummar yvar hällmark u stain.\\
Yvar förfädars bain,\\
yvar martall yvar ain.\\
Yvar maurar u mack,\\
leiksum yvar pinnsvein u rack,\\
skeinar soli igen pa ladingen.\\
\\
Summar de jär när soli kummar yvar Austargarn\\
Summar de jär när brimsar brummar u myggar bleir me ban.\\
Summar de jär när lycku bjaudar sorgi upp till dans.\\
Visst skudd de vare trist skudd de var um summarn inte \\
fanns\\
\\
Summarn kummar...\\
\\
Vackat kan vare mosse pa en gammel sprucken stain,\\
Vackat kan var u hald si bei fast bärgningi är klain.\\
Vackat kan vare sånt som kalles fäult nån annenstans\\
Visst skudd de vare trist skudd de var um Gotland inte \\
fanns.\\
\\
Summarn kummar...\\
\\
{\footnotesize\textit{Text: Allan Nilsson\\ Musik: E Bjuresten}}

\newpage
\visa{DET ÄR VÅREN}
\vspace{10pt}
Det är våren, det är våren,\\
som har kommit igen\\
och med den kärleken.\\
Det är våren, ja det är våren,\\
som har kommit tillbaka igen.\\
När vi varann i portarna krama\\
och kattorna på gårdarna jama,\\
så är det våren, ja det är våren,\\
som har kommit tillbaka igen.\par
\vspace{10pt}
{\footnotesize\textit{Text: Jules Sylvain\\ Musik: Sten Hage}}

\vspace{15pt}
\fvisa{AVSLUTNINGSSÅNGEN}{Den studietid nu varit med tentor och projekt}\\
{\footnotesize\textit{Melodi: Den blomstertid nu kommer}}\\
\\
Den studietid nu varit med tentor och projekt\\
Det varit ganska marigt och några har den knäckt.\\
Ni andra har här samlats och nästan allt är klart.\\
Hur kan nu detta kännas om inte underbart.\\
\\
Nu kommer härlig sommar för oss som än är kvar\\
För er som nu skall knega glöm ej bort hur det var\\
med onsdagar på puben och greven som dessert\\
Ni säkert alltid ångrar att inte stanna här.

\chapterpage{Svenska visor}{Svenska visor}

\chapterpage{Osvenska visor}{Osvenska visor}

\chapterpage{Phulavisor}{Phulavisor}
\fvisa{BAKFYLLOSTROFEN}{Eskimåer jagar valross}
{\footnotesize\textit{Melodi: 34:an}}\\
\\
Eskimåer jagar valross.\\
Alla tyskar jagar älg.\\
Pedofiler jagar småglin.\\
Vilken värld, ja skål och svälj.\\
\\
Uti skogen finns det blåbär.\\
På balkongen står en stol.\\
Uti stolen sitter jag.\\
Jag har vatt bakfull sen i fjol.\\
\\
Om du hör nånting som tickar,\\
kan det kanske va en bomb.\\
Om det är en väckarklocka,\\
bara vänd och somna om.\\
Om du har en liten tax\\
och så en fet gul undulat.\\
\revrpt Ja då hopar sig problemen,\\
man kan inte ge dem samma mat.\rpt\\
\\
{\footnotesize\textit{Text: Johan Tigrelius, David Danowsky,\\ Jonas Hörström, Rickard Nilsson, Annika Svensson}}

\visa{BE- BE- VITAMIN}
{\footnotesize\textit{Melodi: Bä bä vita lamm}}\\
\\
Be-Be- vitamin, finns i brännevin.\\
Mången kalori, simmar däruti.\\
Helgdagssup åt far,\\
och söndagskrök åt mor\\
samt tre små huttar\\
åt lille lille bror.

\fvisa{CLAES PALANDER}{Två snoriga konfirmander}
{\footnotesize\textit{Melodi: Och jungfru hon går i ringen}}\\
\\
Två snoriga konfirmander tog med sig bensin,\\
till Sydsvenskans Claes Palander som skrev om vin.\\
”Se upp jävla gubbe nu åker du dit,\\
för vinet du skrev om det smakade skit”\\
\\
{\footnotesize\textit{Lundakarnevalen 1994}}

\fvisa{INGEMAR BERGMANS SNAPSVISA}{Tystnad"!}
\vspace{10pt}
TYSTNAD!\\
TAGNING!

\visa{KÖTTET KOMMER}
{\footnotesize\textit{Melodi: Vals ur Glada Änkan}}\par
\vspace{10pt}
...Köttet kommer, köttet kommer mört och rött.\\
Gnäggar inte springer inte, det är dött.\\
Skål för alla oxar! Skåla för varje säl.\\
Skål för alla hästar som vi slått ihjäl!\\
Raj raj raj raj raj…\par
\vspace{10pt}
{\footnotesize\textit{Text: Johan Tigrelius, Erik Hedlund,\\ David Danowsky, Linus Åkesson}}

\visa{LAPIN KULTA}
{\footnotesize\textit{Melodi: Broder Jakob}}\\
\\
Lapin Kulta, Lapin Kulta\\
Karjala, Karjala\\
Aura sekä Olvi, Aura sekä Olvi\\
Koff, Koff, Koff\\
Koff, Koff, Koff

\fvisa{ÖL-KANON}{Drick, drick, drick din öl}
{\footnotesize\textit{Melodi: Row, row, row your boat}}\\
\\
Drick, drick, drick din öl\\
låt den rinna ner.\\
Kan du sen kraxa\\
"en laxask med slasktratt”\\
så får du dricka fler.

\fvisa{SI}{W kg m Wb s}
{\footnotesize\textit{Melodi: Studentsången}}\par
\vspace{10pt}
W kg m Wb s\\
$\Omega$ m T A rad\\
cd Sv N s\\
$\Omega$  A m lx dB\\
$^{\circ}$C W/$m^{2}$\\
J/kg H V C\\
kg/$m^{3}$ mol\\
m/$s^{2}$\\
m/$s^{2}$\\
F!\par
\vspace{10pt}
{\footnotesize\textit{Text: Anders Skog}}


\chapterpage{Register}{Register}
\cfoot{}
%\newgeometry{top=1cm,bottom=1.2cm,left=1cm,right=1cm,footskip=0pt,headsep=0.2cm}
\small
\printindex
\hiddenchapterpage{$\emptyset$}
\ohead{}
\newpage
\null
\newpage
\null
\newpage
\null
\newpage
\null
\newpage
\null
\newpage
\null
\newpage
\null
\newpage
\null
\end{document}
